\section{Discussion}

\subsection{Feasibility Achieved, but by a Trivial Equivalence}

The pre-registration framed H1 as the central question: can a richer
detector class (capacity-aware) reach the joint feasibility region
that Paper~11 showed the simple single-$\tau$ family cannot? The
answer is technically yes --- but the H4 rejection reveals that the
detector achieves feasibility by collapsing to the static gate at
both capacity extremes. At K=50 the large $\tau_{50}$ produces zero
fires, so cap\_aware = lag-1 = gated at that capacity. At K=200 the
small $\tau_{200}$ produces 100\% fires, so cap\_aware = fixed = gated
at that capacity.

The substantive contribution is therefore methodological rather
than algorithmic: \textbf{capacity-aware thresholding validates the
static gate as a feasible operating point}. The static gate was
not previously pre-registered as a feasibility candidate; this paper
pins it down by showing that a per-window detector with reasonable
$\tau_K$ choices reduces to it.

\subsection{Why the K=100 Window Fails to Matter}

The intermediate capacity K=100 is the only cell where cap\_aware
differs from the static gate. The Paper~10 default $\tau = 0.05$
fires at 40\% of post-W1 windows; cap\_aware fires the same 40\%
because $\tau_{100} = 0.05$ in the pre-registered vector. cap\_aware
loses $-0.4$~pp to lag-1 at K=100 on cell-mean; the static gate uses
lag-1 unconditionally at K=100 and therefore matches lag-1 exactly.

The K=100 intervention is operationally meaningless: $-0.4$~pp falls
well within the bootstrap noise band for this cell. A
detector-justified procedure at K=100 would need to demonstrate a
material per-window gain, which the data do not provide.

\subsection{Implications for the Detector Family}

The negative H4 result generalises beyond cap\_aware. Any detector
that fires aggressively enough to avoid the K=200 lag-1 hazard
necessarily reverts to fixed at every K=200 window (because the
hazard is per-window and recovers only with a complete revert). The
detector therefore produces the same K=200 decisions as the static
gate. Improvements over the static gate must come from K $\in$ regimes
where the detector fires selectively rather than uniformly --- and
the data in this grid suggest the K=100 regime is too quiet for
selective firing to matter and the K=50 regime is too clear-cut
(lag-1 always helps) for selective firing to be useful.

The H4 rejection is therefore a result about the detector family,
not about $\tau_K = \{50: 0.20, 100: 0.05, 200: 0.02\}$ specifically.
A different $\tau_K$ vector would shift the K=100 firing rate up or
down, but the K=200 and K=50 collapses would remain.

\subsection{Where Richer Detectors Might Help}

The CUSUM and Bayesian detector classes Paper~11 named could in
principle break the K=200 collapse by firing only when cumulative
evidence indicates a sustained shift. If the K=200 W4 window's
$|\Delta| = 0.020$ is a noise excursion that a CUSUM detector would
recognise as low-confidence, the detector could refrain from
reverting to fixed at that window and capture the same $+3.4$~pp
benefit that adaptive05 captured by not firing. This would shift
the cap\_aware K=200 cell mean above the static gate's $0.2664$.

Whether real CUSUM or Bayesian detectors achieve this is a future
work question. The pre-registration discipline that constrained
Papers~7--12 requires that such richer detectors be pre-specified
with locked hyperparameters and decision rules before evaluation.

\subsection{Operational Implication}

The deployable recipe in the studied regime is sharper than Paper~11
concluded:
\begin{itemize}
\item \textbf{Use the static gate.} At $K \leq 100$, deploy lag-1
recalibration; at $K \geq 200$, deploy fixed Paper~4 weights. No
detector, no per-window decision.
\item The capacity-aware detector at the pre-registered $\tau_K$
vector adds operational complexity (a per-window magnitude test)
without producing operationally meaningful gain over the static gate.
\item Adaptive procedures should be preferred only when a richer
detector (CUSUM, Bayesian) has been shown to improve on the static
gate at K=200 in a pre-registered evaluation. The simple
magnitude detector --- whether single-$\tau$ (Paper~10) or
capacity-aware (this paper) --- does not clear that bar in the
studied grid.
\end{itemize}
